\chapter{White piedra}
\index{White piedra}

\section*{Definition and Overview}
White piedra, medically classified as trichomycosis nodularis or tinea nodosa, is a superficial fungal infection of the hair shaft characterized by the formation of soft, whitish, yellowish, or light brown nodules along the hair fibers. The term ``piedra'' originates from the Spanish word for ``stone,'' describing the nodular, bead-like structures that resemble tiny stones clinging to the hair. However, in contrast to the stony-hard and dark nodules of black piedra, the lesions of white piedra are soft, gelatinous, and easily detached. This disease primarily affects the hair shafts of the scalp, pubic area, axilla, beard, mustache, and, in rare instances, eyebrows and eyelashes.

The historical timeline of white piedra dates back to 1865, when the German dermatologist Gustav Beigel first described the condition. He isolated a fungus from infected wigs and human hair, which led to the historical nomenclature of ``Beigel's disease'' and the classification of the causative organism as \textit{Trichosporon beigelii}. Over the subsequent century, advancements in mycological taxonomy, particularly ribosomal RNA sequencing, revealed that the genus \textit{Trichosporon} is highly diverse. Consequently, \textit{T. beigelii} was split into several distinct pathogenic species, including \textit{T. ovoides}, \textit{T. inkin}, and \textit{T. cutaneum}, each exhibiting specific anatomical and epidemiological preferences.

Physiologically, white piedra is a superficial mycosis that is generally restricted to the extra-follicular portion of the hair shaft. The fungus colonizes the cuticle of the hair without invading the hair follicle or the surrounding dermis. In immunocompetent individuals, the condition is benign and represents primarily a cosmetic concern, sometimes accompanied by mild pruritus or structural hair fragility. However, the clinical perspective on the causative organisms has shifted dramatically. Once considered minor dermatological nuisances, \textit{Trichosporon} species are now recognized as significant opportunistic pathogens. In severely immunocompromised patients, such as those with profound neutropenia, hematological malignancies, or advanced HIV infection, these fungi can translocate from superficial colonization sites into the bloodstream, causing disseminated trichosporonosis. This invasive form is characterized by fungemia, multi-organ failure, and high mortality rates, underscoring the clinical importance of identifying and treating superficial infections promptly and effectively.

\section*{Epidemiology}
The epidemiology of white piedra is-intricately linked to environmental factors, geographical location, socioeconomic conditions, and local hygiene practices. The disease is endemic to tropical and subtropical zones, where high ambient temperatures and elevated relative humidity create an ideal environment for the proliferation of yeast-like fungi. Consequently, the highest prevalence is recorded in South America (particularly Brazil, Colombia, and Venezuela), parts of Southeast Asia (such as India, Indonesia, and Thailand), the Middle East, and equatorial regions of Africa.

Despite its association with tropical climates, white piedra is not restricted to these areas. In recent years, an increasing number of sporadic cases have been documented in temperate regions, including the southern United States, parts of Europe, and Japan. This expanding geographical footprint is attributed to global travel, migration, and the creation of localized humid microclimates (such as indoor swimming pools, gyms, or specific working environments).

Demographically, white piedra can affect individuals of any age, gender, or race, although specific sub-populations show higher rates of infection due to behavioral and cultural practices. For instance, in some parts of South America and Asia, scalp white piedra is more commonly observed in young females. This predisposition is linked to long hair and the cultural practice of tying or wrapping wet hair, which traps moisture against the scalp for prolonged periods. Conversely, in other regions, the pubic and axillary variants are more frequently diagnosed in young adult males. Transmission occurs primarily through direct person-to-person contact or indirectly via shared fomites, such as hairbrushes, combs, towels, hats, hair clips, and pillows. Poor personal hygiene, hyperhidrosis (excessive sweating), and bathing in stagnant, contaminated water sources are also significant epidemiological risk factors.

\section*{Aetiology and Pathophysiology}
The development of white piedra involves complex interactions between the causative fungal agent, the hair substrate, and the host's local microenvironment.

\begin{itemize}
    \item \textbf{Causative Organisms}: Fungi of the genus \textit{Trichosporon} are the sole causative agents of white piedra. These are basidiomycetous, anamorphic yeasts that exist in nature as common soil organisms, water saprophytes, and normal flora of the human skin and gastrointestinal tract. Following the taxonomic revision of the genus, several distinct species have been linked to human infections. \textit{Trichosporon ovoides} is the primary agent responsible for white piedra of the scalp, whereas \textit{Trichosporon inkin} is the predominant species isolated from pubic, inguinal, and axillary hair. \textit{Trichosporon cutaneum} is also commonly implicated in both scalp and genital presentations. Other species, such as \textit{Trichosporon asahii} and \textit{Trichosporon mucoides}, are occasionally isolated from hair shafts but are more feared for their ability to cause systemic, deep-seated infections in immunocompromised hosts.
    \item \textbf{Adhesion and Biofilm Formation}: The initial stage of pathogenesis is the attachment of fungal cells (blastoconidia or arthroconidia) to the hydrophobic cuticle of the hair shaft. This attachment is mediated by cell-wall glycoproteins and the secretion of an extracellular mucilaginous polysaccharide matrix. Over time, the proliferating fungal cells organize into a complex, multicellular community resembling a biofilm. This matrix not only anchors the fungus to the hair but also provides a protective barrier against topical antifungal agents and mechanical clearing.
    \item \textbf{Enzymatic Keratin Degradation}: Once firmly attached, \textit{Trichosporon} species synthesize and release a variety of extracellular hydrolytic enzymes, including keratinases, proteases, lipases, and elastases. Keratinases are particularly crucial as they break down the complex keratin proteins that constitute the hair cuticle and cortex. The enzymatic degradation alters the structural integrity of the hair shaft, allowing the fungal hyphae to penetrate beneath the outermost layers of the cuticle.
    \item \textbf{Nodule Maturation}: As the hyphae branch and segment, they undergo arthrosporogenesis, transforming into rectangular-to-oval arthroconidia. Concurrently, blastoconidia multiply by budding. The accumulation of these fungal structures, along with the secreted mucilaginous cement, forms the characteristic nodule. The nodule grows circumferentially, forming a collar around the hair shaft. Because the fungal invasion is mostly ectothrix (external), the hair shaft is initially preserved, but chronic infection leads to complete destruction of the cuticle, exposing the inner cortex and causing the hair to become brittle and break.
    \item \textbf{Host-Microenvironment Synergy}: The growth of \textit{Trichosporon} is heavily promoted by local host factors that increase humidity and temperature. Hyperhidrosis, tight clothing, occlusive hair coverings, and the use of heavy hair oils or grease provide lipid substrates and trap moisture. In healthy hosts, the immune system restricts the infection to the keratinized hair shaft. However, in immunocompromised patients, the lack of functional neutrophils allows the fungus to breach mucosal or cutaneous barriers, leading to angioinvasion and life-threatening hematogenous dissemination.
\end{itemize}

\section*{Clinical Features}
The clinical presentation of white piedra is characterized by localized, non-inflammatory abnormalities of the hair shaft.

\begin{itemize}
    \item \textbf{Characteristics of Nodules}: The pathognomonic sign of white piedra is the presence of multiple, soft, discrete nodules along the infected hair shafts. These nodules typically measure between 1.0 and 1.5 millimetres in diameter, although they can sometimes coalesce to form larger, irregular aggregates. They are soft or gelatinous in consistency, which stands in stark contrast to the hard, grit-like nodules of black piedra.
    \item \textbf{Color and Mobility}: The nodules range in color from white and cream to pale yellow or light brown. They encircle the hair shaft like beads on a string. A key diagnostic feature is their relative mobility: they are not deeply anchored to the hair cortex and can be slid along the length of the hair shaft with gentle finger pressure.
    \item \textbf{Distribution and Localization}: In temperate regions and in female patients, the scalp hair is the most common site of infection. In tropical climates and male patients, there is a higher prevalence of pubic, inguinal, axillary, beard, and mustache hair involvement. Eyebrow and eyelash involvement is rare but can occur, especially in children.
    \item \textbf{Hair Quality and Fragility}: In the early stages of infection, the hair remains structurally sound. However, as the fungal enzymes degrade the cuticle, the hair shaft becomes weak, dry, and brittle. Infected hairs lose their natural sheen and break easily during brushing or washing. This can lead to patchy areas of hair loss (alopecia), though the alopecia is temporary and non-scarring.
    \item \textbf{Subjective Symptoms}: White piedra is typically asymptomatic, and the majority of patients seek medical attention solely for cosmetic reasons. When symptoms do occur, they are generally mild and include localized pruritus, a gritty sensation when touching the hair, or mild irritation of the adjacent skin.
    \item \textbf{Absence of Inflammatory Signs}: The underlying skin, scalp, and hair follicles are normal. There is no erythema, scaling, pustules, or lymphadenopathy. If skin inflammation is present, it suggests a co-existing condition, such as seborrheic dermatitis, or a secondary bacterial infection.
\end{itemize}

\section*{Differential Diagnosis}
Diagnosing white piedra requires differentiation from several other conditions that present with nodules, concretions, or abnormalities on the hair shaft.

\begin{itemize}
    \item \textbf{Black Piedra}: Caused by \textit{Piedraia hortae}, black piedra is characterized by dark, hard, stony nodules that are firmly attached to the hair shaft and cannot be easily slid off. It primarily affects scalp hair and is endemic to tropical rainforests. Microscopic examination of black piedra nodules reveals septate, dark-pigmented hyphae and asci containing ascospores, which are absent in white piedra.
    \item \textbf{Pediculosis Capitis and Pubis (Nits)}: Head and pubic lice lay eggs (nits) that attach to the hair shafts. Nits are teardrop-shaped, translucent or white, and are glued to one side of the hair shaft at an acute angle, rather than encircling it. They do not slide along the hair shaft. Wood's light examination shows a bright yellow-green fluorescence for live nits, and microscopy reveals the classic egg structure with an operculum. The presence of active lice and intense, progressive itching further distinguishes pediculosis.
    \item \textbf{Trichomycosis Axillaris}: Also known as trichobacteriosis, this is a superficial bacterial infection caused by \textit{Corynebacterium} species, most commonly \textit{Corynebacterium flavescens}. It affects axillary and pubic hair, forming yellow, red, or black concretions. These concretions are softer than piedra nodules and are often associated with an unpleasant body odor (bromhidrosis) due to bacterial metabolism of sweat. Gram staining of the concretions reveals Gram-positive coryneform bacteria rather than fungal elements.
    \item \textbf{Trichorrhexis Nodosa}: This is a structural hair shaft defect caused by physical or chemical trauma, such as excessive heat styling, bleaching, or harsh chemical treatments. It presents as small white nodes along the hair shaft. Under a light microscope, these nodes appear as clean fractures of the hair shaft with frayed ends, resembling two paintbrushes pushed against each other. There are no fungal hyphae or spores present.
    \item \textbf{Monilethrix}: A rare, inherited hair disorder characterized by beaded hair shafts. The hair has regular, periodic constrictions (internodes) and swellings (nodes). The nodes are normal hair diameter, while the internodes represent thin, weakened sections where the hair easily fractures. This leads to severe hair fragility and short, stubby hair growth. Unlike white piedra, the beads are part of the hair structure itself and cannot be slid off, and there is no fungal involvement.
    \item \textbf{Peripilous Hair Casts (Pseudonits)}: These are tubular, white-to-yellowish keratinous sheaths that encircle the hair shaft. They represent remnants of the inner and outer root sheaths that have failed to desquamate properly. They are frequently associated with inflammatory scalp disorders like psoriasis, seborrheic dermatitis, or lichen planopilus, or result from traction. While they can slide along the hair shaft like white piedra, microscopy reveals they are composed entirely of keratinous squamous cells without any fungal elements.
\end{itemize}

\section*{When to Seek Medical Care}
Although white piedra is a benign, superficial infection that does not pose an immediate threat to life or limb in healthy individuals, obtaining an accurate diagnosis and treatment plan is essential. Individuals should consult a healthcare provider, such as a general practitioner or a dermatologist, under the following circumstances:
\begin{itemize}
    \item If they notice persistent white, cream, or yellowish nodules on the hair of the scalp, face, or pubic region that do not resolve with standard washing.
    \item If the hair becomes increasingly brittle, split, or begins to break off, leading to areas of hair thinning or localized hair loss.
    \item If the condition is accompanied by moderate-to-severe itching, burning, pain, skin redness, or scaling on the scalp or surrounding skin.
    \item If signs of a secondary bacterial infection develop, such as pus-filled blisters (pustules), increased warmth, swelling, or spreading redness around the hair follicles.
\end{itemize}

For individuals with compromised immune systems, the clinical threshold for seeking medical care is significantly lower. \textit{Trichosporon} species are highly opportunistic and can cause life-threatening invasive infections in vulnerable patients. Immunocompromised individuals (including those undergoing chemotherapy, receiving immunosuppressive therapy for transplants or autoimmune diseases, or living with advanced HIV/AIDS) must seek immediate medical evaluation if they develop any new hair or skin abnormalities.

If an immunocompromised patient experiences signs of systemic infection, such as:
\begin{itemize}
    \item A high fever (above 38 degrees Celsius) or severe chills and shaking.
    \item Rapid heart rate (tachycardia) or rapid breathing.
    \item Unexplained confusion, dizziness, or extreme lethargy.
    \item Shortness of breath or chest pain.
    \item The appearance of painful, necrotic (blackened) skin nodules.
\end{itemize}
They must seek emergency medical services immediately. Call local emergency services in many countries, 911 in the United States and Canada, 999 in the United Kingdom, or proceed to the nearest emergency department without delay.

\section*{Diagnosis and Investigations}
Accurate diagnosis of white piedra is critical to distinguish it from other hair shaft disorders and initiate appropriate therapy. The diagnostic process relies on a combination of clinical inspection and laboratory investigations.

\begin{itemize}
    \item \textbf{Clinical and Physical Assessment}: The clinician begins with a detailed visual and tactile examination of the hair. The presence of soft, light-colored, moveable nodules along the hair shaft, in the absence of scalp inflammation, strongly suggests white piedra.
    \item \textbf{Direct Microscopic Examination (KOH Mount)}: This is the quickest and most definitive diagnostic test. Infected hair shafts are plucked, placed on a microscope slide, and treated with a 10\% to 20\% potassium hydroxide (KOH) solution, sometimes combined with a fungal stain like Parker's ink or calcofluor white. The KOH dissolves the keratin of the hair, allowing clear visualization of the fungal structures. Under light microscopy, the nodule is seen as a dense mass of septate, hyaline hyphae that fragment into rectangular arthroconidia, alongside round-to-oval blastoconidia. The fungal elements primarily encircle the hair shaft (ectothrix invasion).
    \item \textbf{Fungal Culture}: To identify the causative organism, hair samples are inoculated onto culture media. Sabouraud Dextrose Agar (SDA) is standard, but it must be prepared without cycloheximide, as \textit{Trichosporon} species are sensitive to this antibiotic. The cultures are incubated at 25 to 30 degrees Celsius. Within 48 to 72 hours, rapid growth is observed. The colonies are initially white-to-cream, smooth, and yeast-like, later becoming dry, wrinkled, yellowish, and paste-like.
    \item \textbf{Trichoscopy}: Dermoscopy of the hair (trichoscopy) is a non-invasive tool that allows high-magnification visualization of the hair shafts. Under trichoscopy, white piedra appears as irregular, white-to-yellowish, semi-translucent nodules that encircle the hair shaft like a sheath, while the hair follicles and surrounding scalp appear healthy.
    \item \textbf{Biochemical Profiling}: Once a culture is obtained, the specific species of \textit{Trichosporon} can be identified using biochemical systems such as the API 20C AUX or VITEK 2, which test the yeast's ability to assimilate various carbon sources. \textit{Trichosporon} species are also characterized by their ability to produce urease, which distinguishes them from some other yeast genera.
    \item \textbf{Molecular Diagnostics}: In complex, atypical, or systemic cases, molecular methods provide the highest specificity. Polymerase chain reaction (PCR) followed by sequencing of the internal transcribed spacer (ITS) regions of the ribosomal RNA gene allows for rapid, accurate species-level identification.
\end{itemize}

\section*{Management}
Managing white piedra involves a combination of physical grooming, topical antifungal treatments, and, in severe or refractory cases, systemic medications.

\begin{itemize}
    \item \textbf{Physical Hair Removal}: The most rapid and effective intervention is shaving or cutting the hair in the affected areas. Shaving removes the fungal reservoir and prevents further colonization. For infections of the pubic, axillary, or facial hair, shaving is the first-line recommendation. While highly effective, scalp shaving may be unacceptable to some patients due to cosmetic and cultural reasons, necessitating pharmacological management.
    \item \textbf{Topical Antifungal Agents}: Topical therapy is the primary treatment for patients who choose not to shave their hair. Shampoos and creams containing azole antifungals are highly effective. Ketoconazole 2\% shampoo or cream is widely used, applied to the affected area daily and left on for 5 to 10 minutes before rinsing, for a duration of 3 to 6 weeks. Clotrimazole 1\% cream or solution, miconazole 2\% cream, and econazole are effective alternatives.
    \item \textbf{Adjuvant Topical Therapies}: Shampoos containing selenium sulfide 2.5\% or zinc pyrithione 1\% to 2\% can be used in conjunction with azole creams. These agents help reduce fungal load and alter the scalp environment. Topical ciclopirox olamine 1\% is also beneficial due to its broad-spectrum antifungal, antibacterial, and anti-inflammatory activity.
    \item \textbf{Systemic Antifungal Therapy}: Oral antifungal medications are reserved for extensive, recurrent, or treatment-resistant cases, or when a patient refuses scalp hair removal. Oral itraconazole, administered at 100 to 200 mg daily for 3 to 6 weeks, is the gold standard due to its excellent penetration into keratinized tissues and high activity against \textit{Trichosporon} species. Oral fluconazole (100 to 200 mg daily or 150 mg weekly for 4 to 6 weeks) is another option.
    \item \textbf{Hygiene and Sanitization}: To prevent reinfection, patients must sanitize all personal grooming items. Combs, hairbrushes, hairpins, and hairbands should be washed in hot water (above 60 degrees Celsius) or soaked in an antifungal disinfectant. Bed linens, pillowcases, and towels should also be washed at high temperatures. Sharing of these items must be strictly prohibited.
    \item \textbf{Environmental and Behavioral Modifications}: Patients should be advised to dry their hair thoroughly after washing and avoid tying, braiding, or wrapping wet hair, which traps moisture. The use of heavy, occlusive hair oils, pomades, and greases should be suspended, as they provide a lipid-rich environment that supports fungal growth.
\end{itemize}

\section*{Complications}
While white piedra is primarily a localized and benign condition, several complications can arise if it is left untreated or managed inappropriately.

\begin{itemize}
    \item \textbf{Hair Damage and Cosmetic Alopecia}: Over time, the enzymatic action of \textit{Trichosporon} degrades the hair cuticle and invades the cortex. This results in severe weakening of the hair shaft. Hairs become brittle, lose their natural elasticity, and develop split ends. Consequently, mechanical trauma from normal brushing, washing, or styling leads to hair breakage close to the scalp, resulting in patchy, temporary alopecia.
    \item \textbf{Secondary Bacterial Superinfection}: The presence of white piedra nodules, especially in the pubic or axillary regions, can cause mild-to-moderate pruritus. Scratching the affected areas can cause micro-abrasions in the skin. These breaks in the epidermal barrier allow common skin bacteria, such as \textit{Staphylococcus aureus} or \textit{Streptococcus pyogenes}, to enter and infect the tissues, leading to folliculitis, impetigo, or cellulitis.
    \item \textbf{Psychological and Social Impact}: The visible nodules on the hair can be mistaken for head lice (pediculosis) or interpreted as a sign of poor personal hygiene. This misidentification can lead to significant social embarrassment, anxiety, self-consciousness, and social withdrawal, particularly in school-aged children, adolescents, and young adults.
    \item \textbf{Systemic Dissemination (Invasive Trichosporonosis)}: In immunocompromised hosts, superficial colonization by \textit{Trichosporon} species can be a precursor to invasive disease. The fungi can enter the bloodstream through mucosal or cutaneous barriers, leading to disseminated trichosporonosis. This condition is characterized by fungemia, high fever, pulmonary infiltrates, renal involvement, and necrotic skin lesions. Disseminated trichosporonosis is highly lethal, with mortality rates ranging from 50\% to 80\% despite aggressive antifungal therapy.
\end{itemize}

\section*{Prognosis and Outcomes}
The prognosis for individuals with superficial white piedra is excellent, provided there is adherence to the prescribed treatment plan and environmental modifications are made.

When the affected hair is shaved or cut short, the infection is cured immediately, as the fungal substrate is removed. For patients who undergo pharmacological treatment, topical azole therapies or oral antifungal regimens typically achieve complete mycological and clinical clearance within 3 to 6 weeks. The hair structure eventually recovers as new, uninfected hair grows from the follicles, and there is no permanent damage to the hair follicles or scarring of the scalp.

However, recurrence is common if the predisposing factors are not addressed. Failure to sanitize combs and brushes, continued use of occlusive hair products, or persistently wet hair in humid environments can lead to reinfection.

In contrast, the prognosis for invasive, disseminated trichosporonosis in immunocompromised individuals is extremely grave. The causative fungi are often resistant to amphotericin B, the standard treatment for many systemic mycoses, necessitating the use of triazoles (like voriconazole) and recovery of the host's immune system (particularly neutrophil count) for successful clearance. This disparity highlights the critical importance of early diagnosis and prompt eradication of superficial \textit{Trichosporon} colonization in high-risk populations.

\section*{Prevention and Self-Care}
Preventing white piedra and avoiding its recurrence requires a combination of strict personal hygiene, appropriate hair care practices, and environmental awareness.

\begin{itemize}
    \item \textbf{Thorough Drying of Hair}: The most effective preventive measure is ensuring that hair is completely dry after washing. Avoid tying, braiding, or wrapping hair in a towel or headscarf while it is still wet or damp, as this creates a warm, humid microenvironment that encourages fungal spore germination.
    \item \textbf{Hygiene of Grooming Tools}: Do not share personal grooming items such as combs, hairbrushes, hairpins, hairbands, hats, towels, or pillowcases. Wash these items regularly in hot, soapy water, and replace them if an infection is diagnosed.
    \item \textbf{Avoid Occlusive Hair Products}: Limit the use of heavy hair oils, greases, pomades, and thick conditioners, especially in warm and humid climates. These products can trap moisture and organic matter against the hair shaft, providing nutrients for fungal growth.
    \item \textbf{Use of Cotton and Breathable Fabrics}: For prevention of pubic and axillary white piedra, wear loose-fitting undergarments made of natural, breathable fibers like cotton. This helps reduce localized sweating (hyperhidrosis) and moisture accumulation.
    \item \textbf{Proactive Treatment of Sweat Disorders}: If you suffer from excessive sweating (hyperhidrosis) of the scalp or groin, seek medical advice for managing the condition to prevent the creation of a moist environment favorable to fungal colonization.
    \item \textbf{Household Self-Monitoring}: If a family member is diagnosed with white piedra, other household members should inspect their hair shafts for similar nodules. Asymptomatic carriers can act as reservoirs for reinfection.
\end{itemize}

\section*{Sources and Further Reading}
\begin{itemize}
    \item \textbf{DermNet New Zealand}: White Piedra Clinical Reference Guide. Website: \texttt{https://dermnetnz.org/topics/white-piedra}
    \item \textbf{StatPearls (National Center for Biotechnology Information)}: Comprehensive Review on Piedra by J.D. Zuniga and R.M. Mendez. Website: \texttt{https://www.ncbi.nlm.nih.gov/books/NBK540994/}
    \item \textbf{Mycology Online (The University of Adelaide)}: Laboratory Diagnosis and Species Description of the Genus \textit{Trichosporon}. Website: \texttt{https://www.adelaide.edu.au/mycology/}
    \item \textbf{Centers for Disease Control and Prevention (CDC)}: Information on Fungal Diseases and Opportunistic Pathogens. Website: \texttt{https://www.cdc.gov/fungal/}
    \item \textbf{World Health Organization (WHO)}: WHO Fungal Priority Pathogens List to Guide Research, Development, and Public Health Action. Website: \texttt{https://www.who.int/publications/i/item/9789240060241}
    \item \textbf{British Association of Dermatologists (BAD)}: Clinical Guidelines and Resources for Superficial Mycological Infections. Website: \texttt{https://www.bad.org.uk/}
\end{itemize}